\section{Results}\label{sec:results}

All results are from the bench: desktop and laptop computers, a laptop camera,
synthetic frames, loopback and a local network, and for the ARM builds emulation, a
continuous-integration runner and one Raspberry Pi~5. None is from the vehicle.

\subsection{Visibility correction}\label{sec:res-autovis}

On the synthetic scene degraded by Eq.~\eqref{eq:water} (Figure~\ref{fig:autovis}),
after 30 frames the
cast magnitude $m$ fell from 39.1 to 9.1, the 1st--99th percentile luma range rose
from 129 to 249 levels, and the mean absolute difference between neighbouring luma
samples, a measure of visible edge contrast, rose from 7.86 to 19.17. The red patch,
grey-green in the degraded frame ($V - 128 = 0.0$), read red again ($23.1$). On a
uniform frame of luma $30 \pm 3$, the standard deviation rose from 2.00 to 7.80
(3.9$\times$). After a step change in brightness the first frame differed from the
settled result by 105.9 levels in mean luma and the result settled within 1.5~s; a
still scene produced consecutive frames that differ by at most one level in fewer than
0.1\% of pixels. The correction took 0.55~ms per 1280$\times$720 frame on four
threads and 2.09~ms on one (Table~\ref{tab:costs}).

\begin{figure}[!t]
  \centering
  \includegraphics[width=\ifreview 0.8\fi\columnwidth]{figures/fig3_autovis.png}
  \caption{The visibility correction on synthetic frames. Top: the scene degraded by
    Eq.~\eqref{eq:water} (left) and corrected (right). Bottom: the undegraded scene
    (left) and corrected (right), where the correction should change little. The red,
    yellow and orange patches lose most of their red to the model's transmission of
    0.25 and recover only partly.}\label{fig:autovis}
\end{figure}

\subsection{Transport latency}\label{sec:res-latency}

Table~\ref{tab:latency} gives the latency from the frame bus of the sending
\hellion{} to that of the receiving one. Before the two changes on the receiving side
(Section~\ref{sec:transport}) it was 70.4~ms at 640$\times$480 and 30~frames per
second and 37.0~ms at 60: it halved when the frame rate doubled, that is, the chain
held two frame periods. The RTP tap placed both in the decoder's FFmpeg process, one
in the H.264 parser and one in the raw video output; the encoder, from the frame
written to its last RTP packet, took 1.2 to 1.5~ms at 640$\times$480, and
\hellion{}'s own relay, pipe and bus about 0.5~ms. With the delimiter and slice
threading the latency was 3.8~ms at 30 and 3.7~ms at 60~frames per second, no longer
a function of the frame period, and 8.3~ms at 1280$\times$720; the 10th and 90th
percentiles lay within 0.5~ms of the median in every case. A frame whose last packet
is lost still waits for the next frame, since the relay recognizes the end of a frame
by that packet's marker bit.

\begin{table*}[!t]
  \centering\small
  \caption{Transport latency, frame bus to frame bus through two \hellion{} processes
    over loopback; median (10th--90th percentile) of 300 frames.}\label{tab:latency}
  \begin{adjustbox}{max width=\linewidth}
  \begin{tabular}{@{}lccc@{}}
    \toprule
    & \makecell{640$\times$480\\30 frames/s} & \makecell{640$\times$480\\60 frames/s}
    & \makecell{1280$\times$720\\30 frames/s} \\
    \midrule
    Without the delimiter and slice threading & 70.4 (70.2--70.6)~ms & 37.0 (36.8--37.2)~ms & 74.8 (74.6--75.1)~ms \\
    With both & 3.8 (3.7--4.1)~ms & 3.7 (3.6--3.9)~ms & 8.3 (8.1--8.6)~ms \\
    \bottomrule
  \end{tabular}
  \end{adjustbox}
  \tablenote{Without the two measures the decoder held two frame periods: about one in
    FFmpeg's H.264 parser and one in its raw video output. The encoder, from the frame
    written to its last RTP packet, took 1.2--1.5~ms at 640$\times$480, and \hellion{}'s
    own relay, pipe and bus about 0.5~ms. Camera exposure, the tether and the vehicle's
    encoder add to these figures.}
\end{table*}


\subsection{Processing costs}\label{sec:res-costs}

Table~\ref{tab:costs} gives the cost of each stage per 1280$\times$720 frame. On the
desktop processor every stage is a small part of the 33~ms frame period: the
visibility correction 0.55~ms, the blur 0.18 to 1.56~ms over its range of strengths,
the Canny detector 3.4 to 3.8~ms, and publication on the frame bus 34.9~$\mu$s, or
17.2~$\mu$s for a frame of neutral chroma. The measures described in
Section~\ref{sec:operators} reduced the instructions executed by the blur on AArch64
from 900\,M to 283\,M per frame at strength~4 and from 1\,108\,M to 486\,M at
strength~8.

On a Raspberry Pi~5, which has no hardware H.264 encoder, a 1280$\times$720 stream in
placement~B took 0.91 of a core in all, of which the FFmpeg encoder took 0.51; the Pi
was throttled by heat during that run, so both figures are upper estimates. This is
the cost that sending a camera's own H.264 (Section~\ref{sec:transport}) removes from
the vehicle; the saving itself is still to be measured.

\begin{table*}[!t]
  \centering\small
  \caption{Processing costs per 1280$\times$720 frame. x86-64: AMD Ryzen~5 5600,
    minimum over repeated runs, 4 threads unless stated. ARM: executed instructions under
    emulation, one thread, 20 frames.}\label{tab:costs}
  \setlength{\tabcolsep}{4pt}
  \begin{tabular}{@{}L{6.2cm}L{3.2cm}L{4.2cm}@{}}
    \toprule
    Stage & x86-64 & ARM (instructions) \\
    \midrule
    Image statistics and visibility correction & 0.55~ms (2.09~ms, 1 thread) & -- \\
    Blur, strength 1 & 0.18~ms & AArch64 183\,M; ARMv7 263\,M \\
    Blur, strength 4 & 0.59--0.66~ms & AArch64 283\,M; ARMv7 391\,M \\
    Blur, strength 8 & 0.84~ms & AArch64 486\,M; ARMv7 676\,M \\
    Blur, strength 20 & 1.56~ms & -- \\
    Canny, strength 2, thresholds 40/90 & 3.41--3.82~ms & -- \\
    Frame bus publication & 34.9~$\mu$s; 17.2~$\mu$s with neutral chroma & -- \\
    \addlinespace
    \hellion{} preview, 1 and 4 viewers & 4\% and 4\% of a core & -- \\
    Encoder (FFmpeg, libx264), 30 frames/s & 0.21~s of CPU per second of video, laptop & -- \\
    \bottomrule
  \end{tabular}
  \tablenote{Before the measures described in the text, the visibility correction took
    1.49~ms, the blur 0.75, 1.38 and 2.27~ms at strengths 1, 4 and 8 (and 900\,M and
    1\,108\,M instructions on AArch64 at 4 and 8), Canny 3.80--4.20~ms, the preview
    13\% of a core with four viewers, and the encoder 0.30~s of CPU per second of
    video.}
\end{table*}


\subsection{Equivalence and tests}\label{sec:res-equivalence}

Every comparison of Table~\ref{tab:equivalence} was identical: over all cases, thread
counts and the three architectures, the revised code of each change produced the same
bytes as the code it replaced. For a camera's own H.264 sent through the packetizer,
the relay, the decoder and the recorder, every decoded frame and every frame of the
recording was byte-identical to FFmpeg's decode of the source file. The runs between
two machines were functional and gave no measurements.

\begin{table*}[!t]
  \centering\small
  \caption{Equivalence checks: output hashed over all cases for the previous and the
    revised code; every comparison was identical.}\label{tab:equivalence}
  \setlength{\tabcolsep}{4pt}
  \begin{tabular}{@{}L{3.6cm}L{6.2cm}rL{3.2cm}@{}}
    \toprule
    Change & Inputs & Cases & Platforms (threads) \\
    \midrule
    Visibility correction & smooth, flat and noisy synthetic sequences of 60 frames, at
      sizes from 1920$\times$1080 to 1$\times$1, including padded strides & -- & x86-64, AArch64, ARMv7 \\
    Canny & 3 photographs $\times$ 4 strengths $\times$ 7 low $\times$ 6 high thresholds & 504 & x86-64 (1, 3, 4), AArch64, ARMv7 (4) \\
    Blur, half-resolution NEON & 4 sizes $\times$ strengths 0--20 $\times$ 3 or 1 planes & 168 & x86-64 (1, 4), AArch64, ARMv7 (4) \\
    Blur, scratch-plane passes & 7 sizes, including 2-pixel-wide frames $\times$ strengths 0--20 $\times$ 3 or 1 planes & 294 & x86-64 (1, 3, 4), AArch64, ARMv7 (4) \\
    \bottomrule
  \end{tabular}
\end{table*}
